\begin{longtable}{|c|p{4.0cm}|p{2.4cm}|p{2.5cm}|p{3.0cm}|}
\caption{Đối chiếu mức đóng góp đã tuyên bố với bằng chứng thu được}
\label{tab:contribution-levels} \\

\hline
\textbf{Mức} & \textbf{Đóng góp} & \textbf{Khoảng trống tương ứng} & \textbf{Bằng chứng} & \textbf{Đánh giá sau khi có kết quả} \\
\hline
\endfirsthead

\hline
\textbf{Mức} & \textbf{Đóng góp} & \textbf{Khoảng trống tương ứng} & \textbf{Bằng chứng} & \textbf{Đánh giá sau khi có kết quả} \\
\hline
\endhead

\hline
\endfoot

\hline
\endlastfoot

1 & Mô hình ủy quyền nhiều cấp với ba ràng buộc do hợp đồng thi hành
& G1, G3
& Kiểm thử đơn vị hợp đồng, mỗi ràng buộc một phép kiểm khẳng định mã lỗi
& \textbf{Giữ nguyên mức}
\\ \hline

1 & Nguyên tắc bất đối xứng giữa cấp và thu hồi
& G3
& Kịch bản đơn vị hết hạn ủy quyền thu hồi văn bằng đã cấp
& \textbf{Giữ nguyên mức}; phạm vi phát biểu rộng hơn bài toán văn bằng
\\ \hline

1 & Nguyên tắc chuyển quyền gốc phải gồm hai bước
& G2
& Kịch bản chuyển quyền tới địa chỉ không ai kiểm soát
& \textbf{Giữ nguyên mức}
\\ \hline

2 & Kiến trúc gắn kết lớp chứng thực với lớp quản trị trên cùng hạ tầng
& G2
& Một lần truy vấn phục vụ hai bước xác minh
& \textbf{Giữ nguyên mức}
\\ \hline

2 & Quy trình xác minh bảy bước phân biệt loại thất bại
& G4
& Bộ kiểm phát hiện nội dung bị sửa, cả hai cài đặt
& \textbf{Giữ nguyên mức}; kết quả mạnh hơn dự kiến
\\ \hline

3 & Bộ chứng cứ phủ định cho logic phân quyền
& G3, G4
& 83 trên 126 phép kiểm là trường hợp vi phạm có khẳng định mã lỗi
& \textbf{Giữ nguyên mức}
\\ \hline

3 & Phương pháp kiểm chứng tương đương giữa hai cài đặt độc lập
& ---
& 59 phép đối chiếu; ba loại sai lệch được phát hiện
& \textbf{Hạ mức}: chỉ chứng minh tương đương trên tập hành vi được đối chiếu (Mục 4.11.1)
\\ \hline

4 & Ánh xạ mô hình vào bối cảnh giáo dục đại học Việt Nam
& G6
& Phân tích hành lang pháp lý và hạ tầng số hiện có
& \textbf{Giữ nguyên mức}; ở mức phân tích và thiết kế, chưa ở mức tích hợp
\\ \hline

\end{longtable}
